% Terminal note of the homily: everything the speech itself must not carry ---
% the audience, the spoken word count and the pace it implies, the relation to
% the four interpretations of the expansive study, the connections that are
% the homily's own, the translation each quoted text is given in, and the exact
% loci --- followed by the References. Nothing here is read aloud.

\section*{Note on Sources and Delivery}
\runninghead{Sources and Delivery}

{\footnotesize

\textbf{Audience and occasion.} An adult parish assembly, at the Mass
\latin{Dominica prima Adventus}, I classis, of the 1962 \work{Missale
Romanum}, editio typica, pp.~1--2, nos.~1--10, said on Sunday 29~November
2026. The speech presumes that it follows the Gospel, and it speaks of the
Offertory and the Secret as still to come. It names no singer or choir and
holds at a sung or a Low Mass alike: it says that the words of the Gradual and
of the Introit return in the Offertory, not who sings or says them. It uses
the Introit and its repetition at the Offertory, the Collect, the Epistle, the
Gradual, the Alleluia, the Gospel and the Secret; it does not use the
Communion or the Postcommunion. No prayer is recited.

\textbf{Length and pace.} The spoken body runs to about 1,380 words (1,384 by
a plain count with comments and markup removed), which is about 10.6 to 11.5
minutes at an unhurried 120 to 130 words a minute, by arithmetic alone; no one
has delivered or timed it. It was read through in full, silently, for sense
and ease of speech.

\textbf{Relation to the reviewed interpretations.} The homily joins the first
and third of the expansive study's four interpretations. From the first,
\textit{Lift Up Your Heads: An Aging World Ends, and Those Who Do Not Love It
Rejoice}, it takes Gregory's purpose of the prophecy, his aging world, his
reading of the lifted head, his road and his summer, with the Gradual's
waiting set against the Gospel's expectation and the Alleluia heard before the
Gospel. From the third, \textit{The Hour to Rise: The Sleeper Is Told to Wake,
and the One Who Comes Is Asked to Set Him Free}, it takes Chrysostom on Rom
13:11, 12 and 14, Augustine's ``presently'' and his deliverance, the Collect's
petition and the Secret's cleansing. The two are compatible: the third does
not depend on what the signs are, and the first sets Paul's night and
Chrysostom's runners beside Gregory's road and hears them as one call to
haste. Their carrying authors Cyril of Alexandria and Aquinas are left to the
studies, as are the second interpretation (the lights read in the Church), the
fourth (Christ given as man and to be given as God, and the dispute over the
Communion's earth), and Gregory's reading of the calamities of his own day,
the whirlwind of I.1.5--6 among them, together with Augustine's objection to
dating the end from present events (Letter 199), which the speech, using none
of that reading, does not need.

\textbf{What the text changes.} The speech starts from the fear the Gospel
itself describes, men withering ``for fear and expectation'' at the signs, and
the command given in the same breath to the disciples, ``look up and lift up
your heads.'' The decisive textual details are Gregory's \latin{levate capita,
id est exhilarate corda, quia dum finitur mundus, cui amici non estis, prope
fit redemptio quam quaesistis} (I.1.3), which puts the difference between the
two responses in what each loves, and the Gradual's \latin{qui te exspéctant},
before the Gospel, whose verb is that of the Gospel's \latin{exspectatióne}:
two waitings that differ in what is waited for. The Epistle's command to rise,
Augustine's sweet burden of the world and his ``presently'' turn the question
to a hearer who loves the world and knows it; the Collect's \latin{te
liberánte salvári} asks Christ to come and set him free. The fear of the end
is thus understood as a question of love, and answered by a petition the Mass
has already made. With unrelated readings the argument does not survive: each
step rests on Lk 21:25--31, Rom 13:11--14a, Ps 24:1--3, the Alleluia's Ps
84:8, the Collect or the Secret.

\textbf{Connections that are the homily's own.} These are drawn in the
preacher's voice and credited to no Father. Gregory says that grief at the
world's ruin belongs to those who love it (I.1.3); that the Gospel's withering
men and the disciples see the same signs and are divided by what they love
applies his sentence to Lk 21:26 and 28. That the Gradual's \latin{exspéctant}
and the Gospel's \latin{exspectatióne} come from one verb is a fact of the
Missal's Latin; the contrast of the two waitings is the preacher's, as it is
the first interpretation's. That the Alleluia can stand as the prayer of
people who have sought the redemption Gregory says is near joins Ps 84:8 to
his \latin{quaesistis}, as the first interpretation does. That the same haste
as Gregory's \latin{festinare \dots\ quantocius} can be heard in Paul's
command is the preacher's, as it is the first interpretation's. Setting
Augustine's sweet burden of the world after Gregory's love of the world is the
preacher's step. Chrysostom's sentence on v.~12 is told in the first person
plural, as he tells it. That the Collect asks for the deliverance Augustine
named, and that the Church asks Christ to come and set us free so that we can
rise, are the preacher's, as they are the third interpretation's, and so is
the clause that joins the Secret's cleansing to the sins the Collect names;
the Collect asks the Lord to stir up his power and come, and its object is his
power, not the sleeper. The Introit's \latin{levávi} and the Gospel's
\latin{leváte} are one verb; Gregory comments on the Gospel's command and not
on the psalm, and the step by which the psalmist is said to do what that
command asks, and the comparison with the disciples who lift their heads while
the world is shaken, are the preacher's. The response, to lift up to Christ at
the Offertory's ``To thee have I lifted up my soul'' whatever still holds one
while one answers him ``presently'' and to ask him, in the Collect's words, to
set one free from it, applies the antiphon, Augustine's narrative and the
Collect to the hearers. Gregory compares the kingdom to summer (I.1.3);
recalling Christ's command of v.~28 after it is the preacher's close.

\textbf{Sources, translations and loci.} Each source is listed once, with the
translation the speech uses; ``paraphrased'' means told in the speech's own
words. \emph{The Mass}, nos.~1--10. The Collect is quoted in the Lasance
English (p.~103) from ``that, defended by Thee''; its petition to stir up his
might and come is paraphrased. The Secret's ``cleansing us by their mighty
power'' is quoted from Lasance (p.~105), and the rest of the prayer is
paraphrased in its words. The Lasance \work{New Roman Missal} (1945) is a
historical study aid and not an approved liturgical translation of the 1962
Missal. \emph{Scripture} is quoted in the Douay--Rheims, Challoner revision:
Lk 21:26 and 28; Rom 13:11--12a, with ``passed, and'' for the Gutenberg text's
``passed And'', and 13:14a; Ps 24:1b--2a for the Introit, without the Douay's
``O Lord'', which the Missal's antiphon lacks; Ps 24:3b for the Gradual,
without its opening ``For'', as Lasance gives the Gradual (p.~104); Ps 84:8
for the Alleluia; Jas 4:4b. Lk 21:20, 24, 25, 29--31, Rom 13:12b and Ps 24:3a
are told in the speech's own words; the Douay's ``wait on thee'' renders the
Vulgate's \latin{sustinent te}, where the Missal's chant reads \latin{te
exspéctant}. \emph{St Gregory the Great}, in the expansive study's English of
the Migne Latin, in places rendered afresh from the Latin: \work{Homiliae in
Evangelia} I.1.1, \latin{paratos nos invenire desiderans, senescentem mundum
quae mala sequantur denuntiat, ut nos ab ejus amore compescat}, with ``to
restrain us from loving it'' quoted, and \latin{ut si Deum metuere in
tranquillitate non volumus, vicinum ejus judicium vel percussionibus attriti
timeamus}, paraphrased closely; I.1.5, the aging body and the world weighed
down by its age, paraphrased, and \latin{Nolite ergo, fratres mei, diligere
quem videtis diu stare non posse}, quoted; I.1.3, \latin{levate capita, id est
exhilarate corda}, paraphrased, with ``because while the world of which you
are not friends is ending, the redemption you have sought draws near'' quoted;
\latin{caput pro mente} and \latin{Levare itaque capita est mentes nostras ad
gaudia patriae coelestis erigere}, paraphrased; \latin{Qui ergo Deum diligunt
ex mundi fine gaudere atque hilarescere jubentur, quia videlicet eum quem
amant mox inveniunt, dum transit is quem non amaverunt}, paraphrased closely;
his citation of Jas 4:4 and \latin{radices cordis in ejus amore plantaverunt},
paraphrased; \latin{festinare ad ea quantocius debemus \dots\ breviore via
pervenire}, paraphrased; \latin{Quid est vita mortalis, nisi via? Et quale
sit, fratres mei, perpendite, in labore viae lassescere, et tamen eamdem viam
nolle finiri}, quoted; \latin{fructus mundi ruina est. Ad hoc enim crescit, ut
cadat}, paraphrased, and \latin{quia tunc moeroris nostri nubila transeunt, et
vitae dies aeterni solis claritate fulgescunt}, quoted. The inscription says
the homily was preached to the people in St Peter's basilica; the speech says
``to the people of Rome''. \emph{St John Chrysostom}, in the Nicene and
Post-Nicene English: \work{Homilies on Romans} 24, on 13:11, earnestness that
would ``wither down to nothing'' and the racers who ``rouse themselves up the
more'' near the end of the course, paraphrased; on 13:12, the neighbours woken
before dawn, ``It is day now!'', and the dressing and leaving of dreams,
quoted and paraphrased; on 13:14, works, arms and ``the Lord Himself for a
garment,'' with Christ ``laid round about us as a garment'' and the soul ``a
dwelling for Himself,'' paraphrased. \emph{St Augustine}, \work{Confessions}
VIII.5.12, in Pilkington's English (NPNF): ``with the baggage of the world was
I sweetly burdened, as when in slumber,'' paraphrased; the man who desires to
awake and is steeped again in drowsiness, and who ``even after it is time to
rise with pleasure yields to it,'' paraphrased; ``Awake, thou that sleepest,''
``Presently, lo, presently,'' ``Leave me a little while,'' ``But `presently,
presently,' had no present,'' ``who shall deliver me from the body of this
death'' and ``Thy grace only, through Jesus Christ our Lord,'' quoted; ``the
violence of custom, whereby the mind is drawn and held, even against its
will,'' paraphrased; VIII.12.29, his reading of Rom 13:13--14, told.

\textbf{Text checked.} The Lasance English was taken from
\texttt{propers/verified.md}, and the Douay--Rheims from the registered
Challoner verse texts of Luke, Romans, Psalms and James, whose hashes were
recomputed and matched. Gregory was re-read in the registered Wikisource Latin
text (hash recomputed and matched) at I.1.1--6; Chrysostom in the registered
NPNF passages of \work{Homilies on Romans} 24 on 13:11, 12 and 14 and their
NPNF vol.~11 text (hash matched); the \work{Confessions} in the registered
NPNF vol.~1 text (hash matched) at VIII.5.12 and VIII.12.29.

\par}

\Needspace{14\baselineskip}
\section*{References}
\runninghead{References}

{\footnotesize
\begin{itemize}\setlength{\itemsep}{0.15em}\setlength{\parskip}{0pt}
\item \work{Missale Romanum ex decreto Sacrosancti Concilii Tridentini restitutum, Summorum Pontificum cura recognitum}, editio typica (Typis Polyglottis Vaticanis, 1962), CMAA facsimile, Proprium de Tempore, pp.~1--2, nos.~1--10.
\item Lasance, F.~X., and Walsh, F.~A., \work{The New Roman Missal in Latin and English}, revised edition (New York: Benziger Brothers, 1945), ``First Sunday of Advent,'' pp.~103--105 (Collect, Gradual, Secret).
\item \work{The Holy Bible Translated from the Latin Vulgate} (Douay--Rheims, Challoner revision), Project Gutenberg text: Ps 24:1--3; 84:8; Lk 21:20--31; Rom 13:11--14; Jas 4:4.
\item St Augustine, \work{Confessions} VIII.5.12 and VIII.12.29, trans.\ J.~G. Pilkington, Nicene and Post-Nicene Fathers, 1st ser., vol.~1.
\item St John Chrysostom, \work{Homilies on Romans} 24, on 13:11, 12 and 14, Nicene and Post-Nicene Fathers, 1st ser., vol.~11.
\item St Gregory the Great, \work{Homiliae in Evangelia} I.1, nn.~1, 3 and 5, Migne Latin text (Wikisource).
\end{itemize}
\par}
